%!TEX root = ../graduate-work.tex
\phantomsection
\section*{Аннотация}
\addcontentsline{toc}{section}{Аннотация}

Формальная верификация нейронных сетей~--- математическое доказательство
того, что сеть удовлетворяет заданным свойствам безопасности для всех
входных данных из допустимой области,~--- сталкивается с фундаментальным
ограничением: алгоритмы распространения границ (IBP, CROWN,
$\alpha$-CROWN) требуют хранения матриц весов целиком на одном GPU,
что делает невозможной верификацию моделей с миллиардами параметров.

В данной работе исследуется применение методов параллелизма, разработанных
для обучения больших языковых моделей, к задаче формальной верификации.
Реализованы и экспериментально исследованы два подхода в рамках
фреймворка \texttt{auto\_LiRPA} / \texttt{$\alpha,\beta$-CROWN}:

\textbf{Tensor Parallelism (TP)}~--- шардирование матриц весов и
матриц линейной релаксации по нескольким GPU. Достигнуто $\approx\!2\times$
снижение пиковой памяти при $P=2$ GPU; soundness bounds подтверждена
на бенчмарках VNN-COMP 2022 (MNIST-FC).

\textbf{Fully Sharded Data Parallelism (FSDP)}~--- послойное шардирование
только матриц весов с AllGather перед каждым слоем. Bounds побитово
идентичны однопроцессорному режиму: экономия baseline-памяти составила
50\%, пиковой~--- 34--39\% на широких MLP. FSDP интегрирован
с полной верификацией ($\beta$-CROWN + Branch-and-Bound) и расширен
на свёрточные слои (\texttt{BoundConv}); корректность верифицирована
на CIFAR-100 ResNet medium/large (VNN-COMP 2024), для ResNet-large
получен полный результат верификации (unsat).

Экспериментально установлено, что в режиме $\alpha$-CROWN+BaB узким
местом по памяти являются alpha-тензоры (per-neuron slopes),
а не матрицы весов. Определены границы применимости каждого подхода
и сформулированы направления дальнейшей работы.
